% =====================================================================
\section{Своя веб-страница: HTML, CSS и JavaScript}\label{les:web}
% =====================================================================
\lessonintro
  {написать страницу управления вручную и понять, что происходит внутри готовых панелей вроде ESPUI}
  {урок 5 (Wi-Fi и ESPUI), урок 3 (кнопка), урок 4 (ШИМ), урок 8 (общее состояние под мьютексом)}
  {страница с кнопкой и ползунком, написанная на HTML, CSS и JavaScript, управляет светодиодом}
  {Wi-Fi-сеть на \SI{2.4}{GHz}; светодиод с резистором на GPIO4}

До сих пор панели управления мы собирали из готовых элементов ESPUI: библиотека сама отдавала
браузеру страницу. Этот урок --- про то, \emph{как} такие страницы устроены. Вы напишете страницу
сами и научите плату отвечать на её запросы. Материал сложнее предыдущих: здесь сразу три языка
(HTML, CSS, JavaScript) и две программы (на плате и в браузере). Не торопитесь и проходите примеры по порядку.

\subsection{Первый веб-сервер}

\codecap{Первый веб-сервер}
\codeinput{l08_hello_web}

Загрузи программу и подключи телефон к сети \code{S3-Web} (пароль \code{lesson-44-demo}) --- плата
раздаёт её сама, как в уроке~5, роутер не нужен. Открой в браузере адрес из Serial Monitor
(обычно \code{http://192.168.4.1/}). Обнови страницу несколько раз --- число секунд будет расти:
страницу каждый раз заново «готовит» плата.

\begin{itemize}
  \item \code{WiFi.softAP(имя, пароль)} --- плата создаёт свою сеть;
  \item \code{WebServer server(80)} --- создать сервер. 80 --- номер «двери» (порта), в которую браузер
        стучится по умолчанию;
  \item \code{server.on("/", handleRoot)} --- «если спросят адрес \code{/}, вызови \code{handleRoot}»;
  \item \code{server.send(200, тип, текст)} --- отправить ответ. 200 --- код «всё хорошо»;
  \item \code{server.handleClient()} --- проверить, не пришёл ли запрос. Вызываем постоянно в \code{loop()}.
\end{itemize}


\subsubsection{Пример: управление через GET}

Самый простой способ управлять светодиодом --- ссылки. Нажатие на ссылку \code{/toggle} переключает
светодиод, а форма с ползунком отправляет \code{/level?v=...}:

\codecap{Светодиод управляется GET-запросами}
\codeinput{l08_get_led}

\begin{figure}[H]
\centering
\begin{tikzpicture}[y=-0.85cm]
  \node[blk,fill=blkmem,minimum width=3cm] (b) at (0,0) {Браузер};
  \node[blk,fill=blkcpu,minimum width=3cm] (s) at (9,0) {ESP32-S3};
  \draw[thick,gray] (b.south) -- (0,8.2);
  \draw[thick,gray] (s.south) -- (9,8.2);
  \draw[arr,accent] (0,1.5) -- node[above,font=\ttfamily\scriptsize]{GET /} (9,1.5);
  \draw[arr,okgreen] (9,2.6) -- node[above,font=\ttfamily\scriptsize]{200 OK + HTML-страница} (0,2.6);
  \node[font=\sffamily\scriptsize,anchor=east] at (-0.2,3.7) {нажали «Вкл/выкл»};
  \draw[arr,accent] (0,4.4) -- node[above,font=\ttfamily\scriptsize]{GET /toggle} (9,4.4);
  \fill[esp!30] (8.85,4.6) rectangle (9.15,5.6);
  \node[font=\ttfamily\scriptsize,anchor=west] at (9.3,5.1) {ledOn = !ledOn};
  \draw[arr,okgreen] (9,6.0) -- node[above,font=\ttfamily\scriptsize]{303 → откройте /} (0,6.0);
  \draw[arr,accent] (0,7.0) -- node[above,font=\ttfamily\scriptsize]{GET /} (9,7.0);
  \draw[arr,okgreen] (9,7.9) -- node[above,font=\ttfamily\scriptsize]{200 OK (Светодиод: ВКЛ)} (0,7.9);
\end{tikzpicture}
\caption{Что происходит при нажатии кнопки на странице}
\end{figure}

Работает, но у этого способа два минуса: команды меняют состояние через GET (см. совет выше), а
страница каждый раз полностью перезагружается и мигает. Исправим оба.

\subsubsection{Пример: отправка данных через POST}

Сделаем «записку для платы»: форма отправляет текст методом POST, плата печатает его в Serial Monitor
и показывает на странице.

\codecap{Приём данных методом POST}
\codeinput{l08_post_note}

В форме указано \code{method='POST'}, а на плате обработчик зарегистрирован с \code{HTTP\_POST}:
на обычный GET-запрос по адресу \code{/note} он не ответит. Текст записки достаётся функцией
\code{server.arg("text")} --- так же, как параметры GET, только данные пришли из тела запроса.

\subsection{HTML: из чего сделана страница}

Страница, которую отправляет сервер, написана на языке \textbf{HTML} (англ. \emph{HyperText Markup Language} ---
язык разметки гипертекста). Это не язык программирования: он описывает, \emph{что} есть на странице.
Каждый элемент обозначается \textbf{тегами} в угловых скобках: открывающим \code{<p>} и закрывающим \code{</p>}.
Теги вкладываются друг в друга, как матрёшки:

\begin{figure}[H]
\centering
\begin{tikzpicture}[
  lvl/.style={draw,thick,rounded corners=3pt,font=\ttfamily\small,anchor=north west,inner sep=4pt}]
  \node[lvl,fill=blkmem!40,minimum width=75mm,minimum height=52mm] (html) at (0,0) {};
  \node[font=\ttfamily\small,anchor=north west] at (0.1,-0.05) {<html>};
  \node[lvl,fill=blkrtc!40,minimum width=71mm,minimum height=9mm] at (0.2,-0.55) {};
  \node[font=\ttfamily\small,anchor=north west] at (0.3,-0.6) {<head> \textcolor{gray}{название, стили}};
  \node[lvl,fill=blkrf!50,minimum width=71mm,minimum height=36mm] at (0.2,-1.6) {};
  \node[font=\ttfamily\small,anchor=north west] at (0.3,-1.65) {<body> \textcolor{gray}{то, что видно}};
  \node[lvl,fill=white,minimum width=67mm] at (0.4,-2.2) {<h1>Мой светодиод</h1>};
  \node[lvl,fill=white,minimum width=67mm] at (0.4,-2.95) {<p>Яркость: 50 \%</p>};
  \node[lvl,fill=white,minimum width=67mm] at (0.4,-3.7) {<button>Вкл / выкл</button>};
  \node[lvl,fill=white,minimum width=67mm] at (0.4,-4.45) {<input type="range">};
  % результат
  \begin{scope}[xshift=9.2cm]
    \fill[black!80,rounded corners=4mm] (0,-5.3) rectangle (3.4,0.2);
    \fill[white] (0.2,-4.9) rectangle (3.2,-0.3);
    \node[font=\bfseries\small,anchor=west] at (0.3,-0.8) {Мой светодиод};
    \node[font=\small,anchor=west] at (0.3,-1.7) {Яркость: 50 \%};
    \node[draw,fill=gray!15,font=\scriptsize,anchor=west] at (0.3,-2.5) {Вкл / выкл};
    \draw[gray,thick] (0.3,-3.3) -- (3.0,-3.3);
    \fill[accent] (1.65,-3.3) circle (0.12);
    \node[font=\sffamily\scriptsize,text=white] at (1.7,-5.05) {так без CSS};
  \end{scope}
\end{tikzpicture}
\caption{Структура HTML-страницы (слева) и как браузер её показывает (справа)}
\end{figure}

\begin{table}[H]
\centering\small
\begin{tabular}{@{}ll@{}}
\toprule
Тег & Что означает\\
\midrule
\code{<h1>}\ldots\code{</h1>} & крупный заголовок (есть ещё \code{h2}--\code{h6} поменьше)\\
\code{<p>}\ldots\code{</p>} & абзац текста\\
\code{<b>}\ldots\code{</b>} & жирный текст\\
\code{<a href="/toggle">}\ldots\code{</a>} & ссылка\\
\code{<button>}\ldots\code{</button>} & кнопка\\
\code{<input type="range">} & ползунок (\code{type="text"} --- поле для текста)\\
\code{<form>}\ldots\code{</form>} & форма: собирает поля и отправляет на сервер\\
\code{<div>}\ldots\code{</div>} & «коробка» для группировки элементов\\
\bottomrule
\end{tabular}
\caption{Теги, которых хватит для первых страниц}
\end{table}

Теперь напишем такую страницу сами. Раньше мы собирали HTML из кусков строк, и в нём легко запутаться
в кавычках. Удобнее записать всю страницу одной строкой в форме \code{R"rawliteral(} \ldots
\code{)rawliteral"}: внутри можно писать кавычки и переносы строк как есть. Откройте страницу, затем
измените текст, добавьте ещё один пункт списка и снова загрузите скетч.

\codecap{Страница на HTML без оформления}
\codeinput{l08_html}

\subsection{CSS: делаем красиво}

HTML отвечает за содержание, а за \emph{внешний вид} --- цвета, шрифты, отступы, скругления, тени ---
отвечает язык \textbf{CSS} (англ. \emph{Cascading Style Sheets} --- каскадные таблицы стилей). Стили пишут
внутри тега \code{<style>}. Каждое правило состоит из \emph{селектора} (к чему применить) и списка свойств:

\begin{figure}[H]
\centering
\begin{tikzpicture}
  \node[font=\ttfamily\large,inner sep=0] (c) {\textcolor{esp}{button} \{ \textcolor{accent}{background}: \textcolor{okgreen!60!black}{\#1f618d}; \textcolor{accent}{border-radius}: \textcolor{okgreen!60!black}{12px}; \}};
  \draw[esp,thick] ($(c.north west)+(0.5,0.1)$) -- ++(0,0.4) node[above,font=\sffamily\scriptsize,text=black]{селектор: все кнопки};
  \draw[accent,thick] ($(c.north west)+(3.3,0.1)$) -- ++(0,0.4) node[above,font=\sffamily\scriptsize,text=black]{свойство};
  \draw[okgreen!60!black,thick] ($(c.north west)+(5.6,0.1)$) -- ++(0,0.4) node[above,font=\sffamily\scriptsize,text=black]{значение};
  \node[font=\sffamily\scriptsize,below=2mm of c] {читается так: «у всех кнопок фон цвета \#1f618d и скруглённые на 12 пикселей углы»};
\end{tikzpicture}
\caption{Устройство правила CSS}
\end{figure}

Цвета в CSS часто записывают шестнадцатеричным кодом \code{\#RRGGBB}: две цифры на красный, зелёный
и синий --- та же смесь трёх цветов, что в RGB-светодиоде из урока~\ref{les:led}! Например, \code{\#ff0000} ---
красный, \code{\#ffffff} --- белый.

\begin{table}[H]
\centering\small
\begin{tabular}{@{}ll@{}}
\toprule
Свойство CSS & Что меняет\\
\midrule
\code{color}, \code{background} & цвет текста и фона\\
\code{font-family}, \code{font-size} & шрифт и его размер\\
\code{padding}, \code{margin} & внутренний и внешний отступ\\
\code{border-radius} & скругление углов (\code{50\%} --- круг)\\
\code{box-shadow} & тень или свечение вокруг элемента\\
\code{width}, \code{height} & ширина и высота\\
\code{transition} & плавная смена стиля (анимация)\\
\bottomrule
\end{tabular}
\caption{Свойства CSS для первых экспериментов}
\end{table}

\subsection{Практика: панель управления светодиодом}

Соберём всё вместе: HTML-страницу с CSS-стилями и маленькой программой на \textbf{JavaScript} --- языке,
который выполняется прямо в браузере. Он будет отправлять POST-запросы \emph{без перезагрузки страницы}
и перерисовывать «лампочку» на экране.

\begin{figure}[H]
\centering
\definecolor{webbg}{HTML}{EEF2F7}\definecolor{webred}{HTML}{FF5A4F}\definecolor{webblue}{HTML}{1F618D}
\begin{tikzpicture}
  % телефон
  \fill[black!85,rounded corners=5mm] (0,0) rectangle (4.4,8);
  \fill[webbg] (0.25,0.6) rectangle (4.15,7.4);
  \fill[black!60,rounded corners=1mm] (1.7,7.6) rectangle (2.7,7.75);
  % карточка
  \fill[black!10,rounded corners=3mm] (0.62,1.28) rectangle (3.82,6.88);
  \fill[white,rounded corners=3mm] (0.55,1.35) rectangle (3.75,6.95);
  \node[font=\sffamily\bfseries,text=espdark] at (2.15,6.45) {Мой светодиод};
  \fill[webred,opacity=0.25] (2.15,5.2) circle (0.85);
  \fill[webred] (2.15,5.2) circle (0.62);
  \fill[webblue,rounded corners=2mm] (1.1,3.55) rectangle (3.2,4.1);
  \node[font=\sffamily\small,text=white] at (2.15,3.82) {Вкл / выкл};
  \draw[gray!60,line width=2pt,line cap=round] (0.9,2.85) -- (3.4,2.85);
  \draw[accent,line width=2pt,line cap=round] (0.9,2.85) -- (2.4,2.85);
  \fill[accent] (2.4,2.85) circle (0.13);
  \node[font=\sffamily\small] at (2.15,2.1) {Яркость: \textbf{60} \%};
  % выноски
  \tikzset{co/.style={font=\sffamily\scriptsize,align=left,anchor=west},cl/.style={thin,gray}}
  \draw[cl] (3.75,6.6) -- (5.2,6.9) node[co]{\code{.card} --- белая карточка:\\ скругление и тень};
  \draw[cl] (2.8,5.2) -- (5.2,5.4) node[co]{\code{.led.on} --- круг со\\ свечением \code{box-shadow}};
  \draw[cl] (3.2,3.8) -- (5.2,3.9) node[co]{\code{button} → POST \code{toggle}};
  \draw[cl] (3.4,2.85) -- (5.2,2.6) node[co]{\code{input type=range} →\\ POST \code{level=60}};
  \draw[cl] (3.2,2.1) -- (5.2,1.5) node[co]{JavaScript вписывает\\ число из ответа JSON};
\end{tikzpicture}
\caption{Как будет выглядеть панель на телефоне}
\end{figure}

Собираем панель по шагам.

\subsubsection{Шаг 1. Плата умеет отвечать JSON}

Сначала сделаем только «серверную» часть, без страницы. Обработчик \code{/api} на запрос POST меняет
состояние, а на любой запрос отвечает текущим состоянием в формате \textbf{JSON}: это текст вида
\code{\{"on":true,"level":60\}}, который одинаково легко читают и люди, и программы. Откройте
\code{/api} в браузере --- вы увидите ответ платы.

\codecap{Шаг 1. Адрес /api отвечает JSON}
\codeinput{l08_panel_s1}

\subsubsection{Шаг 2. Страница с кнопкой и ползунком}

Теперь добавим страницу: HTML-элементы (заголовок, кнопка, ползунок) и маленькую программу на JavaScript.
Функция \code{send()} отправляет POST-запрос командой \code{fetch()}, а \code{show()} вписывает ответ
платы в страницу. Оформления пока нет --- страница некрасивая, но уже работает.

\codecap{Шаг 2. HTML и JavaScript без оформления}
\codeinput{l08_panel_s2}

\subsubsection{Шаг 3. Оформление на CSS}

Остаётся сделать красиво. Блок \code{<style>} вставляется в \code{<head>}; в тело страницы добавляются
«карточка» \code{<div class="card">} вокруг элементов и круг \code{<div id="led" class="led">} --- «лампочка»,
которую скрипт перекрашивает. Полный скетч с оформлением лежит в папке \code{sketches} (\code{l08\_panel}).

\codecap{Шаг 3. Оформление: блок style (фрагмент)}
\codeinput{l08_panel_css}

\subsubsection{Как это работает}

\begin{enumerate}
  \item При открытии адреса платы браузер делает \code{GET /} и получает страницу \code{PAGE}.
  \item Скрипт на странице сразу делает \code{GET /api}. Плата отвечает состоянием в формате
        \textbf{JSON} --- \code{\{"on":true,"level":60\}}. Это текст, который легко читают и программы, и люди.
  \item При нажатии кнопки функция \code{fetch()} отправляет \code{POST /api} с телом \code{toggle=},
        а при движении ползунка --- \code{level=60}. Плата меняет яркость и снова отвечает JSON.
  \item Функция \code{show()} перекрашивает «лампочку» и обновляет число. Страница не перезагружается.
\end{enumerate}

\begin{figure}[H]
\centering
\begin{tikzpicture}[y=-0.8cm]
  \node[blk,fill=blkmem,minimum width=3cm] (b) at (0,0) {Браузер + JavaScript};
  \node[blk,fill=blkcpu,minimum width=3cm] (s) at (9,0) {ESP32-S3};
  \draw[thick,gray] (b.south) -- (0,6.3);
  \draw[thick,gray] (s.south) -- (9,6.3);
  \draw[arr,accent] (0,1.3) -- node[above,font=\ttfamily\scriptsize]{GET /} (9,1.3);
  \draw[arr,okgreen] (9,2.2) -- node[above,font=\ttfamily\scriptsize]{200 OK + страница с CSS и JS} (0,2.2);
  \draw[arr,accent] (0,3.1) -- node[above,font=\ttfamily\scriptsize]{GET /api} (9,3.1);
  \draw[arr,okgreen] (9,4.0) -- node[above,font=\ttfamily\scriptsize]{\{"on":false,"level":50\}} (0,4.0);
  \node[font=\sffamily\scriptsize,anchor=east] at (-0.2,4.7) {двинули ползунок};
  \draw[arr,accent] (0,5.0) -- node[above,font=\ttfamily\scriptsize]{POST /api \quad level=60} (9,5.0);
  \fill[esp!30] (8.85,5.1) rectangle (9.15,5.6);
  \draw[arr,okgreen] (9,5.9) -- node[above,font=\ttfamily\scriptsize]{\{"on":true,"level":60\}} (0,5.9);
\end{tikzpicture}
\caption{Страница загружается один раз, дальше идут только маленькие запросы к \code{/api}}
\end{figure}

\begin{tip}
Хранить страницу в строке \code{PAGE} удобно, пока она небольшая. Для больших сайтов с картинками
файлы записывают в файловую систему платы (LittleFS); работа с ней разобрана в уроке 18.
\end{tip}


\subsection{Самостоятельный опыт: кнопка и браузер вместе}

Добавьте в скетч кнопку из урока 3: нажатие и ссылка \code{/toggle} должны менять одно и то же
состояние. Если кнопку читает отдельная задача (урок 8), защитите \code{ledOn} и \code{level}
мьютексом --- обработчики WebServer выполняются из loop(), а задача кнопки может работать параллельно.

\begin{trouble}
\begin{itemize}[leftmargin=*]
  \item \textbf{Страница не открывается} --- телефон в другой (например, гостевой) сети или адрес набран с ошибкой.
  \item \textbf{Страница открывается очень долго} --- в \code{loop()} есть долгие \code{delay()}, и \code{handleClient()} вызывается редко.
  \item \textbf{Кнопка на странице ничего не делает} --- откройте консоль браузера (F12): там видно ошибки JavaScript и коды ответа платы.
\end{itemize}
\end{trouble}

\begin{summary}
HTML описывает содержание, CSS задаёт оформление, JavaScript отправляет запросы и обновляет страницу.
GET читает состояние, POST передаёт команду; JSON связывает браузер с программой платы.

\begin{itemize}[leftmargin=*]
  \item \code{WebServer}: адрес → обработчик; \code{handleClient()} вызываем постоянно.
  \item Сеть --- ещё один источник команд; устройство обрабатывает их теми же функциями, что и нажатие кнопки.
  \item ESPUI делает ту же работу за вас: генерирует страницу и обмен данными по WebSocket.
\end{itemize}
\end{summary}

\begin{selfcheck}
\begin{enumerate}
  \item Что происходит между браузером и платой при переходе по ссылке \code{/toggle}?
  \item Чем отличаются GET и POST? Почему POST сам по себе не шифрует данные?
  \item За что отвечают HTML, CSS и JavaScript?
  \item Какое CSS-свойство превращает квадрат в круг?
\end{enumerate}
\end{selfcheck}

\begin{tasks}
\begin{enumerate}
\item \lv{1} Подключитесь к сети, выведите IP-адрес и откройте панель. Найдите в коде HTML, CSS и JavaScript.
\item \lv{1} Измените заголовок, подписи и цвета панели. Сделайте индикатор зелёным; объясните, какие CSS-селекторы изменили.
\item \lv{2} Добавьте кнопки яркости 30, 60 и 100 процентов. Каждая отправляет POST и отображает ответ платы.
\item \lv{2} Добавьте GET /api/uptime с временем работы в JSON и обновляйте число на странице каждые две секунды без перезагрузки.
\item \lv{3} Если за 15 секунд подключение не удалось, запустите точку доступа и ту же панель. Проверьте оба сценария.
\item \lv{4} Добавьте физическую кнопку и синхронизацию двух браузеров: периодически читайте состояние, показывайте ошибку сети и восстанавливайте связь без повторного переключения светодиода.
\end{enumerate}
\end{tasks}
